% s01-prelim.tex —— 印刷页 2–4（§0 预备）
\section*{0.\ 预备}
\addcontentsline{toc}{section}{0.\ 预备}
\srcpage{2}

本文讨论处处非分歧的情形。此时人们希望对每个三元组建立 Drinfeld–Langlands 对应。三元组由底域 $k$、系数域 $C$，以及 $k$ 上有限型概形范畴的 $C$-系数上同调理论组成。

三个主要的三元组（底域、系数域、上同调理论）是：
\begin{itemize}
\item[(Betti)] $k=\C$，$C$ 是特征 $0$ 的代数闭域\yizhu{原文系数域用斜体 $C$ 记，而 Betti 情形的底域取 $k=\C$（复数域）；两者字母相同而含义不同，勿混。}{术语背景，I（原文 p2；上下文推断）}，以及经典拓扑下 $C$-向量空间的可构造层论，
\item[(De Rham)] $k$ 的特征为 $0$，$C=k$，以及 Zariski 拓扑下 holonomic $D$-模的理论，
\item[($\ell$-进)] $k$ 任意，$C=\Ql$，其中素数 $\ell$ 在 $k$ 中可逆，以及 étale 拓扑下 $\ell$-进层的理论。
\end{itemize}

Drinfeld–Langlands 对应的首要数据，是底域 $k$ 上的一条代数曲线。对此曲线，人们附上一些模空间，它们一般是 $k$ 上的代数叠。因此我将自由使用代数叠的语言（参 [L-M]）。

上述上同调理论对代数叠范畴尚未得到系统的建立，文献也零散。对 Frenkel、Gaitsgory 与 Vilonen 在其文章中所用的 $\ell$-进三元组，情况相当令人满意。事实上，他们为获得几何 Fourier 变换还假设 $k$ 的特征为 $p>0$，但这一限制并不必要。Beilinson 与 Drinfeld 的工作 [B-D] 使用 De Rham 三元组，其中可以找到代数叠上微分算子环的定义。Betti 三元组原则上最基本，Bernstein 与 Luntz [B-L] 研究过它，但解析叠的理论尚待写出。

本文仍然使用最后这个 Betti 三元组。因此底域为 $k=\C$。设 $\cS$ 是有限型代数叠，或更一般地是 $\C$ 上局部有限型的代数叠。对与 $\cS$ 相伴的解析叠上具有有界可构造上同调的 $C$-向量空间层复形，其导出范畴我简记为 $\Dbc(\cS)$。当 $\cS$ 表示为概形 $S$ 按代数群 $G$ 的作用所得的商时，该解析叠正是 $S(\C)^{\mathrm{an}}$ 按 $G(\C)^{\mathrm{an}}$ 的作用所得的商叠\yizhu{原文用同一个字母 $S$ 既记代数叠又记用来呈现它的概形，此处按上下文区分。}{术语背景，I（原文 p2；上下文推断）}。此时 $\Dbc(\cS)$ 就是 Bernstein 与 Luntz 在 [B-L] 中引入的导出范畴 $D^{b}_{G(\C)^{\mathrm{an}}}(S(\C)^{\mathrm{an}})$。

对本文考虑的所有代数叠 $\cS$，导出范畴 $\Dbc(\cS)$ 带有张量积与对偶函子 $D:\Dbc(\cS)^{\mathrm{opp}}\to\Dbc(\cS)$。对这类叠之间所考虑的可表态射 $f:\cS\to\cT$，\srcpage{3}有前推函子 $f_*,f_!:\Dbc(\cS)\to\Dbc(\cT)$ 与拉回函子 $f^*,f^!:\Dbc(\cT)\to\Dbc(\cS)$，它们满足 Grothendieck 六种运算的形式体系。范畴 $\Dbc(\cS)$ 带有中间反常 $t$-结构，其心是 $\cS$ 上反常层的范畴 $\Perv(\cS)$（参 [B-B-D]）。回忆：$\Dbc(\cS)$ 的这个满子范畴是阿贝尔的、Noether 的、Artin 的，并且关于对偶 $D$ 自对偶。

设 $\pi:\cV\to\cS$ 是常秩 $r$ 的向量丛。可考虑 $\cV$ 按乘法群 $\Gm$ 的数乘作用所得的 $\cS$-商叠 $\overline{\cV}$。若 $\cV^\circ$ 是 $\cV$ 中零截面的补开集，则包含 $\cV^\circ\subset\cV$ 诱导开浸入
\[ j:\Pv(\cV)=[\cV^\circ/\Gm]=\overline{\cV}{}^\circ\hookrightarrow\overline{\cV}, \]
其中 $\Pv(\cV)\to\cS$ 是 $\cV\to\cS$ 的直线所成的射影丛。该开浸入的余闭子叠，是 $\cV$ 的零截面按 $\Gm$ 的平凡作用所得的商，即分类叠 $B(\Gm/\cS)$。

$\cV$ 上 $\Gm$-等变的 $C$-向量空间可构造复形，按定义就是 $\Dbc(\overline{\cV})$ 的对象。

对这些 $\Gm$-等变复形，存在几何 Fourier 变换，称为与向量丛 $\cV/\cS$ 相伴的齐次 Fourier 变换（参 [Lau 3]）。这个 Fourier 变换是导出范畴的等价
\[ \Four_{\overline{\cV}/\cS}:\Dbc(\overline{\cV})\to\Dbc(\overline{\cV}^\vee), \]
其逆为 $\Four_{\overline{\cV}^\vee/\cS}$。这里 $\overline{\cV}^\vee$ 是对偶向量丛 $\pi^\vee:\cV^\vee\to\cS$ 按乘法群 $\Gm$ 的数乘作用所得的 $\cS$-商叠。它与对偶交换，并且是 $t$-正合的。可以如下定义它。

在商叠 $\cA=[\A^1/\Gm]$ 上，即在仿射直线按乘法群的数乘作用所得的商上，有复形 $\Psi=\beta_*C\in\operatorname{ob}\Dbc(\cA)$。这里 $\beta:\Spec(\C)=[\Gm/\Gm]\hookrightarrow\cA$ 是由包含 $\Gm\subset\A^1$ 诱导的开浸入。此复形的非零上同调层为
\[ \cH^0(\Psi)=C \]
以及
\[ \cH^1(\Psi)=\alpha_*C, \]
其中 $\alpha:B(\Gm)\hookrightarrow\cA$ 是与 $\beta$ 互补的闭浸入，由原点到 $\A^1$ 的包含诱导。实际上 $\Psi$ 不是不可约反常层，它是 $\alpha_*C[-1]$ 通过点状常数反常层 $C$ 的扩张。

向量丛 $\cV$ 与其对偶丛之间的自然配对态射，过商后成为代数叠的态射
\[ \mu:\overline{\cV}^\vee\times_\cS\overline{\cV}\to\cA. \]
\srcpage{4}%
又若记 $p:\overline{\cV}^\vee\times_\cS\overline{\cV}\to\overline{\cV}$ 与 $p^\vee:\overline{\cV}^\vee\times_\cS\overline{\cV}\to\overline{\cV}^\vee$ 为两个典范投影，则按定义
\[ \tag*{(0.1)} \Four_{\overline{\cV}/\cS}(K)=(p^\vee)_!\bigl(p^*K\otimes\mu^*\Psi\bigr)[r-1],\qquad \forall K\in\Dbc(\overline{\cV}). \]

态射 $p$ 与 $p^\vee$ 不是可表态射，且 $(p^\vee)_!$ 不保持 $\Dbc$：它把 $\Dbc(\overline{\cV}^\vee\times_\cS\overline{\cV})$ 送到 $D^-_c(\overline{\cV}^\vee)$ 中。然而可以验证，$\Four_{\overline{\cV}/\cS}$ 本身保持 $\Dbc$。

\begin{definitionZ}[0.2]
设 $L$ 是 $\Pv(\cV)$ 上的反常层。若遗忘支撑的态射 $j_!L\to j_*L$ 是同构，则称 $L$ 是 proper 的\yizhu{此处的 proper 是 (0.2) 给出的反常层术语，与态射的「紧合」（propre）无关，故保留原文不译。}{术语背景，A（原文 (0.2)；参 BBD，Astérisque 100）}。此时 $j_!L\cong j_*L$ 实际上是 $\overline{\cV}$ 上的反常层，它正是 $L$ 的中间延拓 $j_{!*}L$。
\end{definitionZ}

\begin{lemmaZ}[0.3]
设 $L$ 是 $\Pv(\cV)$ 上的反常层。则 $L$ 是 proper 的，当且仅当 $(\bar\pi^\circ)_!L=0$，其中 $\bar\pi^\circ:\Pv(\cV)\to\cS$ 是典范投影。\hfill$\square$
\end{lemmaZ}

与射影丛 $\Pv(\cV)\to\cS$ 相伴的几何 Radon 变换，是函子
\[ \Rad_{\Pv(\cV)/\cS}:\Dbc(\Pv(\cV))\to\Dbc(\Pv(\cV^\vee)), \]
它由
\[ \Rad_{\Pv(\cV)/\cS}(L)=(q^\vee)_!q^*L[r-1] \]
定义。这里 $q:\cH\to\Pv(\cV)$ 与 $q^\vee:\cH\to\Pv(\cV^\vee)$ 是关联叠
\[ \cH\subset\Pv(\cV^\vee)\times_\cS\Pv(\cV) \]
的两个典范投影。该关联叠由 $\cV\to\cS$ 的一条直线与一个超平面构成的对组成，其中直线含于该超平面。关于这个变换的详细研究，参看 Brylinski 的专著 [Br]。

\begin{lemmaZ}[0.4]
设 $L$ 是 $\Pv(\cV)$ 上的反常层，并假设它是 proper 的。则 $\Four_{\overline{\cV}/\cS}(j_!L)$ 也是 proper 的。它在开集 $\Pv(\cV^\vee)=\overline{\cV}^{\vee\circ}$ 上的限制等于 $\Rad_{\Pv(\cV)/\cS}(L)$；该开集是 $\overline{\cV}^\vee$ 中零截面的补集。\hfill$\square$
\end{lemmaZ}
